\documentclass[10pt,a4paper]{amsart}
\usepackage[margin=19mm]{geometry}
\usepackage{amsmath,amssymb,mathtools}
\usepackage[scr=rsfs,cal=euler]{mathalfa}
\usepackage{xfrac}
\usepackage[unicode,hidelinks,bookmarks=false]{hyperref}
\newcommand{\ie}{i.e.\ }
\newcommand{\cf}{cf.\ }
\newcommand{\resp}{respectively\ }
\newcommand{\Subsection}{\subsection}
\providecommand{\nobreakdash}{\nobreak\hbox{-}}
\setlength{\emergencystretch}{2em}
\multlinegap0pt
\usepackage{algorithm}\usepackage{algpseudocode}
\usepackage{tikz}
\usepackage{enumerate}
\usepackage{float}
\usepackage{amssymb}
\newtheorem{theorem}{Theorem}
\newtheorem{lemma}{Lemma}
\newtheorem{claim}{Claim}
\newcommand{\bP}{\mathbb{P}}
\newcommand{\cH}{\mathcal{H}}
\title{\bf\Large Most probably trangle-free graphs}
\author{Yuhang Bai, Gyula O.H. Katona, Zixuan Yang}
\date{Source: arXiv:2602.22782v1 (CC BY 4.0)}
\begin{document}
\maketitle

\noindent\emph{Source and licence.} \bf\Large Most probably trangle-free graphs. Version arXiv:2602.22782v1 (CC BY 4.0), distributed by arXiv under the Creative Commons Attribution 4.0 International licence, http://creativecommons.org/licenses/by/4.0/, as recorded in the arXiv licence metadata for this exact version and shown on the abstract page https://arxiv.org/abs/2602.22782v1. Full licence notice: \texttt{licenses/arxiv-2602.22782-CC-BY-4.0.txt}.\par\smallskip

\noindent\textit{Editorial scope.} Counted unit: the article's lemma lem:linear3 (source0.tex lines 251--257). The article's complete original proof of it is delivered with it (source0.tex lines 259--346, one \texttt{proof} environment of 88 lines). No mathematics is written, repaired or re-derived: the statement and the proof are the article's own lines, in the article's own notation, and no retained line is substituted or re-flowed.  Retained uncounted as the source-local context the proof needs: lines 237--243.  Nothing was omitted from any retained span, so the package carries no omission record.  The retained lines cite 1 works, whose entries are transcribed verbatim from the article's own bibliography source in the e-print and delivered with short editorial labels recorded in the item's reference mapping.  No retained line contains a figure environment, an image-inclusion command or drawing code, so no image is delivered and the package declares no asset dependency.  Editorial formatting only, in addition: an AMS \texttt{amsart} preamble that restates the article's own declarations and macros used by the retained lines, namely the environments \texttt{theorem}, \texttt{lemma}, \texttt{claim} on separate counters, together with the standard packages those lines need.  The retained environments print in this document as Theorem 1, Lemma 1, Claim 1 in order of appearance, whereas the article prints the same items with the numbering of its own full document.  The complete native source of the article as distributed in the arXiv e-print for this exact version is bundled unchanged as \texttt{sources/195352/source0.tex}.
\begin{theorem}[Harris, \cite{Har}]\label{Har-ineq}
    Let $A_1, A_2,\ldots, A_m \subseteq \Omega$ be events in the independent Bernoulli space $\{0,1\}^n$.
    If $A_1,A_2,\ldots,A_m$  are all decreasing, then
    \[
    \mathbb{P}(\bigcap_{i\in [m]}A_i ) \geq \prod_{i\in [m]}\mathbb{P}(A_i).
    \]
\end{theorem}

\begin{lemma}\label{lem:linear3}
Let $\cH$ be a linear $3$-uniform hypergraph with vertex set $V(\cH)$ and edge set $E(\cH)$ where $|V(\cH)| = n$ and $|E(\cH)| = r$. Let $S\subseteq V(\cH)$
be formed by including each vertex independently with probability $p\in(0,1)$. Then
$$
\mathbb P(S \text{ is an independent set in } \cH)\le 1-p+p(1-p^2)^r.
$$
\end{lemma}

\begin{proof}
We prove the result by using induction on the number of edges $r$.
% For  the case of $r=1$, let $e=\{x,y,z\}$ be the edge of $\cH$. It is clear that every subset $S\subseteq V(\cH)$ which does not contain at least one of $\{x,y,z\}$ is an independent set, and the corresponding probability equals $1-p^3$. Thus the result holds. 
For the case of $r = 0$, It is clear that the result holds. 

Next we may assume that $r\ge 1$, and the statement holds for $r-1$. 
 Let $\cH'= \cH -\{e_0\}$, where  $e_0=\{x_0,y_0,z_0\}$ is an edge in $ \cH$.
We  define two events as follow.
\begin{itemize}
\item Event $A$: the subset $S$  is an independent set in $\cH'$.
\item Event $B$: the subset $S$ contains the vertices $x_0,y_0,$ and $z_0$.
\end{itemize}
Note that $S$ is an independent set in $\cH$ if and only if the event $A$ occurs and the event $B$ does not occur. So we have
\begin{align}\label{S-ind}
\mathbb P(S \text{ is an independent set in } \cH)&=\mathbb P(A\cap \overline{B})\notag\\
&=\mathbb P(A)-\mathbb P(A\cap B)\notag\\
&=\mathbb P(A)-p^3\,\mathbb P(A\mid B).
\end{align}

\begin{claim}\label{cl:condtiona}
$\mathbb{P}(A\mid B)\ge (1-p^2)^{r-1}$.
\end{claim}

\begin{proof}
We identify each subset $S\subseteq V(H)=\{v_1,\dots,v_n\}$ with its characteristic vector
$\omega_S\in\Omega=\{0,1\}^n$, where
\[
\omega_S(v_i)=
\begin{cases}
1 & \text{if } v_i\in S,\\
0 & \text{if } v_i\notin S.
\end{cases}
\]
There is a partial ordering on the vectors of $\Omega$:  $\omega\le \omega'$ if for all $v_i$ we have 
$\omega(v_i)\le \omega'(v_i)$. On the level of sets this  means that $\omega_S\le \omega_{S'}$ if and only if $S\subseteq S'$. 
For each edge $e\in E(\cH')$, we define
\begin{itemize}
\item event $A_e$: the subset $S$  does not contain all three vertices of $e$.
\end{itemize}
For  $S'\subseteq S$, since $\omega_S \in A_e$   implies that $\omega_{S'} \in A_e$, it follows from the definition that $A_e$ is a decreasing event. Then by the definition of event $A$, we have 
\begin{align}\label{AAe}
A=\bigcap_{e\in E(\cH')}A_e.
\end{align}

We assert that $\mathbb P(A_e\mid B)\ge 1-p^2$ for all $e\in E(\cH')$. Since $\cH$ is linear, every edge $e \in E\left(\cH^{\prime}\right)$ satisfies $\left|e \cap e_0\right| \leq 1$, where $e_0=\{x_0,y_0,z_0\}$.  If $|e \cap e_0|=0$,  then we have 
$$
\mathbb{P}(\overline{A_e} \mid B)=p^3,
$$
and thus 
\begin{align*}
\mathbb{P}\left(A_e \mid B\right)=1-p^3 > 1-p^2,
\end{align*}
as desired. So we consider the case when $\left|e \cap e_0\right|=1$. Without loss of  generality, let $x_0$ be the common vertex of $e\cap e_0$. The event $B$ implies that  $x_0 \in S$, and the other two vertices of $e$ are still included independently with probability $p$. This implies that
\begin{align*}
\mathbb{P}(\overline{A_e} \mid B)=p^2, 
\end{align*}
and thus 
\begin{align*}
 \mathbb{P}\left(A_e \mid B\right)=1-p^2,
 \end{align*}
the assertion holds.

Condition on the event $B$, which by definition forces $x_0,y_0,z_0\in S$.
Under this conditioning, the inclusion of each vertex $v\in V\setminus\{x_0,y_0,z_0\}$ is still decided
independently with probability $p$. 
Thus, under the conditional distribution $\bP(\,\cdot\,\mid B)$, we may apply Harris's inequality (i.e., Theorem \ref{Har-ineq}) and equality~(\ref{AAe}) to derive 
\begin{align*}
\mathbb P(A\mid B)&=\mathbb P\Big(\bigcap_{e\in E(\mathcal H')}A_e \Big| B\Big)\\
&\ge \prod_{e\in E(\mathcal H')}\mathbb P(A_e\mid B)
\ge (1-p^2)^{r-1},
\end{align*}
which completes the proof of the claim.
\end{proof}


The induction hypothesis applied to $\cH'$ yields
$$
\mathbb P(A)\le 1-p+p(1-p^2)^{r-1}.
$$
Thus, by Claim~\ref{cl:condtiona} and (\ref{S-ind}),  we have 
\begin{align*}
\mathbb P(S \text{ is an independent set  in } \cH)&=\mathbb P(A)-p^3\,\mathbb P(A\mid B)\\
&\le \mathbb P(A)-p^3(1-p^2)^{r-1}\\
&\le 1-p+p(1-p^2)^{r-1}-p^3(1-p^2)^{r-1}\\
&=1-p+p(1-p^2)^r.
\end{align*}
This completes the proof of Lemma~\ref{lem:linear3}.
\end{proof}

\begin{thebibliography}{99}
\bibitem[Har]{Har} T. Harris, A lower bound for the critical probability in a certain percolation process, \emph{Math. Proc. Cambridge Philos. Soc. }, \textbf{1} (1960), 13-20.
\end{thebibliography}
\end{document}
